% TOP_PROBLEM: {"id": 3350, "title": "Odd perfect numbers", "category_id": 1, "category": {"id": 1, "name": "number_theory", "display_name": "Number Theory", "description": "Properties of integers, prime numbers, Diophantine equations.", "slug": "number-theory", "order_index": 1, "created_at": "2026-07-31T15:26:25.670Z"}, "status": "open", "problem_number": "OPG-2147", "set_id": 12}
% FIRST_POSTED: 2026-10-01
% ENTRY_KIND: statement_only
% UnsolvedMath ID 3350; source catalogue code OPG-2147
% !TeX program = lualatex
% Definition and sources only; this file is not an LLM solution attempt.
\documentclass[11pt,a4paper]{article}
\usepackage[margin=25mm]{geometry}
\usepackage{fontspec}
\setmainfont{lmroman10-regular.otf}[BoldFont=lmroman10-bold.otf,
 ItalicFont=lmroman10-italic.otf,BoldItalicFont=lmroman10-bolditalic.otf]
\usepackage{amsmath,amssymb,mathtools}
\usepackage{unicode-math}
\setmathfont{latinmodern-math.otf}
\AtBeginDocument{\let\setminus\smallsetminus}
\usepackage{xurl,hyperref}
\hypersetup{colorlinks=true,urlcolor=blue}
\setcounter{secnumdepth}{0}
\setlength{\parindent}{0pt}
\setlength{\parskip}{5pt}
\setlength{\emergencystretch}{3em}
\begin{document}
\section*{Odd perfect numbers}\label{3350}
{\small UnsolvedMath ID 3350 · OPG-2147 · Number Theory · OpenGarden}
\subsection{Definitions and mathematical statement}
For a positive integer $n$, define
$\sigma(n)=\sum_{d\mid n,\ d>0}d$, the sum of all positive divisors,
including $1$ and $n$. A perfect number satisfies $\sigma(n)=2n$;
equivalently, its proper positive divisors sum to $n$. The conjecture
asserts that every perfect number is even:
\[
 \forall n\in\mathbb Z_{\ge1},\qquad
 n\equiv1\pmod2\ \Longrightarrow\ \sigma(n)\ne2n.
\]
For a completely explicit arithmetic formulation, write an odd integer
as $n=\prod_{i=1}^r p_i^{a_i}$ with distinct odd primes $p_i$ and
positive integer exponents $a_i$. Multiplicativity of the divisor sum
then gives the equivalent nonexistence problem
\[
 \prod_{i=1}^r(1+p_i+\cdots+p_i^{a_i})
       =2\prod_{i=1}^r p_i^{a_i}.
\]
The quantifier includes every number $r$ of prime factors and every
choice of exponents; imposing any fixed bound on these parameters is
only a restricted version. The case $n=1$ is not perfect since
$\sigma(1)=1$. The required conclusion is either a proof of
nonexistence throughout the odd positive integers or an explicit odd
integer with a verified divisor-sum equality.
\subsection{Short English statement}
Is it impossible for the proper divisors of an odd positive integer to add up to that integer?
\subsection{Sources}
\begin{itemize}
\item[S1] ULAM.AI and the UnsolvedMath Contributors, supplied problems.json, record 3350 (OPG-2147), OpenGarden collection. Catalogue identity and source transcription; CC BY 4.0.
\url{https://huggingface.co/datasets/ulamai/UnsolvedMath}.
\item[S2] Open Problem Garden, Odd perfect numbers. Original problem and discussion; the mathematical formulation below is expanded from the supplied source record.
\url{http://www.openproblemgarden.org/op/odd_perfect_number}.
\end{itemize}
\end{document}
